% =====================================================================
%  style.tex — feuille de style LaTeX (pipeline lualatex)
%  Injectée dans le préambule par build.sh (--include-in-header).
%  Source de vérité du rendu : cadre, quadrillage de fond, polices,
%  titres, blocs de code, citations, sommaire.
%  Palette : or #C8880A / brun #5C3D00 (cohérente avec l'ancien style.css).
% =====================================================================

% ---- Couleurs ----
\definecolor{or}{HTML}{C8880A}
\definecolor{brun}{HTML}{5C3D00}
\definecolor{brunfonce}{HTML}{7A5200}
\definecolor{creme}{HTML}{FDF7EE}
\definecolor{cremecode}{HTML}{F9F5EE}
\definecolor{bordcode}{HTML}{E0C882}
\definecolor{grille}{HTML}{F0E9D8}   % quadrillage de fond, très pâle
\definecolor{shadecolor}{HTML}{F9F5EE} % fond des blocs de code (pandoc Shaded)

% ---- Polices (lualatex / fontspec) + repli emoji couleur ----
% La fonction de repli rend les emoji (🙂😳…) via Noto Color Emoji.
\usepackage{luaotfload}
\directlua{luaotfload.add_fallback("emojirepli", {"NotoColorEmoji:mode=harf;"})}
\setmainfont{P052}[RawFeature={fallback=emojirepli}]            % Palatino (corps)
\setsansfont{Nimbus Sans}[RawFeature={fallback=emojirepli}]     % titres
\setmonofont{DejaVu Sans Mono}[Scale=0.85]                      % code

% ---- Interligne un peu aéré ----
\linespread{1.12}

% =====================================================================
%  Fond de page : quadrillage léger + cadre simple + numéro de page
%  Dessiné sur CHAQUE page via eso-pic, sauf quand \pageframefalse
%  (utilisé pour la couverture pleine page).
% =====================================================================
\usepackage{eso-pic}
\usepackage{tikz}

\newif\ifpageframe \pageframetrue

% Marges (doivent rester synchronisées avec celles passées à pandoc)
\newcommand{\margeHaut}{22mm}
\newcommand{\margeBas}{25mm}
\newcommand{\margeGauche}{22mm}
\newcommand{\margeDroite}{22mm}

\AddToShipoutPictureBG{%
  \ifpageframe
    \begin{tikzpicture}[remember picture, overlay]
      % Quadrillage très pâle, limité à la zone de texte
      \begin{scope}
        \clip ([shift={(\margeGauche,\margeBas)}]current page.south west)
           rectangle ([shift={(-\margeDroite,-\margeHaut)}]current page.north east);
        \draw[step=5mm, line width=0.15pt, grille]
          (current page.south west) grid (current page.north east);
      \end{scope}
      % Cadre simple doré
      \draw[or, line width=0.7pt, rounded corners=1pt]
        ([shift={(14mm,14mm)}]current page.south west)
        rectangle ([shift={(-14mm,-14mm)}]current page.north east);
      % Numéro de page (centré, entre le cadre et le texte)
      \node[anchor=south] at ([shift={(0,16.5mm)}]current page.south)
        {\sffamily\small\color{brun}\thepage};
    \end{tikzpicture}%
  \fi
}
\pagestyle{empty} % pas de numéro natif : géré par le fond ci-dessus
% Neutralise le style « plain » (utilisé par \tableofcontents) pour éviter
% un second numéro de page superposé à celui du cadre.
\makeatletter
\renewcommand{\ps@plain}{\ps@empty}
\makeatother

% =====================================================================
%  Titres
% =====================================================================
\usepackage{titlesec}

% ## (chapitres pandoc = \subsection) : grands, sans-serif brun, filet doré,
% chacun commençant sur une nouvelle page.
\titleformat{\subsection}[hang]
  {\sffamily\bfseries\LARGE\color{brun}}{}{0pt}{}
  [{\vspace{3pt}{\color{or}\titlerule[1.4pt]}}]
\titlespacing*{\subsection}{0pt}{4pt}{14pt}

% ### (sous-parties pandoc = \subsubsection) : sans-serif or, taille moyenne.
\titleformat{\subsubsection}[hang]
  {\sffamily\bfseries\large\color{or}}{}{0pt}{}
\titlespacing*{\subsubsection}{0pt}{16pt}{5pt}

% Nouvelle page avant chaque chapitre (## -> \subsection)
\let\origsubsection\subsection
\renewcommand{\subsection}{\clearpage\origsubsection}

% =====================================================================
%  Citations (blockquote)
% =====================================================================
\usepackage{tcolorbox}
\tcbuselibrary{breakable, skins}
\renewenvironment{quote}
  {\begin{tcolorbox}[breakable, enhanced, colback=creme, boxrule=0pt,
     leftrule=4pt, colframe=or, arc=0pt, left=14pt, right=10pt,
     top=8pt, bottom=8pt, before skip=10pt, after skip=10pt]\itshape\color{black!75}}
  {\end{tcolorbox}}

% =====================================================================
%  Blocs de code
% =====================================================================
% pandoc enveloppe le code coloré dans un environnement Shaded.
% On le redéfinit en boîte crème opaque + fin cadre doré : le fond masque
% proprement le quadrillage et le code se détache du texte.
\usepackage{etoolbox}
\AtBeginDocument{%
  \renewenvironment{Shaded}
    {\begin{tcolorbox}[breakable, enhanced, colback=cremecode,
       colframe=bordcode, boxrule=0.5pt, arc=1.5pt,
       left=8pt, right=8pt, top=6pt, bottom=6pt,
       before skip=9pt, after skip=9pt]}
    {\end{tcolorbox}}%
}
% Renvoi à la ligne des longues lignes de code (sinon elles débordent du cadre).
% La flèche ↪ marque visuellement une ligne coupée.
\usepackage{fvextra}   % fournit breaklines pour les blocs Verbatim/Highlighting
\fvset{breaklines=true, breakanywhere=true,
       breaksymbolleft={\raisebox{0.6ex}{\small\color{or}$\hookrightarrow$}}}

% =====================================================================
%  Tables : un peu d'air
% =====================================================================
\renewcommand{\arraystretch}{1.25}

% =====================================================================
%  Sommaire (\tableofcontents injecté à la place du HTML)
% =====================================================================
\usepackage{tocloft}
\renewcommand{\cfttoctitlefont}{\sffamily\bfseries\LARGE\color{brun}}
\setcounter{tocdepth}{2}          % n'affiche que les ## (chapitres)
% Entrées de chapitre : alignées à gauche, sobres, points de conduite dorés
\renewcommand{\cftsubsecfont}{\normalfont}
\renewcommand{\cftsubsecpagefont}{\bfseries\color{or}}
\renewcommand{\cftsubsecindent}{0pt}
\renewcommand{\cftsubsecnumwidth}{0pt}
\renewcommand{\cftsubsecleader}{\color{grille}\cftdotfill{\cftdotsep}}
\renewcommand{\cftdotsep}{2}
\setlength{\cftbeforesubsecskip}{6pt}

% =====================================================================
%  Listes : un peu plus compactes
% =====================================================================
\usepackage{enumitem}
\setlist{itemsep=2pt, topsep=4pt, parsep=0pt}
